\section{Calibration on the triangular cell}\label{cal:section}\label{sec:calibration}

We construct the radial functions used by the Voronoi test potential.
The potential has the same gradient as the triangular Green function
on a regular hexagon, while its energy on every other sector reduces
to two scalar estimates.  The continuation outside that hexagon is
auxiliary; its explicit form makes those estimates accessible to
interval arithmetic.

\subsection{The local triangular Green function}
Set
\begin{equation}\label{cal:geometry}
 a=\sqrt{\frac{2}{\sqrt3}},\qquad p=\frac a2,\qquad
 R_0=\frac{2p}{\sqrt3},\qquad x_0=\frac{R_0}{2}.
\end{equation}
Thus $p$ and $R_0$ are the inradius and circumradius of the covolume-one
triangular Voronoi hexagon, and $x_0=p\tan(\pi/6)$.
Identify the plane with $\mathbb C$, and write
$\partial_z=(\partial_x-i\partial_y)/2$, with a nearest lattice point on the
positive real axis.  For $m\ge1$ put
\begin{equation}\label{cal:coefficients}
 k_m=6m,\qquad
 c_m=\frac1{k_m}\sum_{z\in\Lambda_\triangle\setminus\{0\}}z^{-k_m}.
\end{equation}
These sums converge absolutely and are real by reflection symmetry.
We also introduce an auxiliary index $0$, with $k_0=6$.
It is distinct from index $1$: both exponents equal $6$, but they will
carry different radial functions.

\begin{lemma}\label{cal:green-expansion}
The mean-zero triangular Green function has the expansion
\begin{equation}\label{cal:green-series}
 h_\triangle(z)=-\log|z|+C_\triangle+\frac{\pi|z|^2}{2}
                 +\sum_{m\ge1}c_m\operatorname{Re}(z^{6m}),
 \qquad 0<|z|<a,
\end{equation}
where $C_\triangle$ is its regular part at the origin.  The series and
all of its derivatives converge locally away from $|z|=a$.
\end{lemma}
\begin{proof}
Away from the lattice, $\Delta h_\triangle=2\pi$, so
$2\partial_z^2h_\triangle$ is a meromorphic elliptic function.  Its double
pole at each lattice point has principal part $(z-\lambda)^{-2}$.
Subtract the absolutely locally convergent Weierstrass series
\[
 \wp(z)=\frac1{z^2}+\sum_{\lambda\ne0}
       \bigl((z-\lambda)^{-2}-\lambda^{-2}\bigr).
\]
The difference is holomorphic and periodic, hence constant.  If
$\omega=e^{\pi i/3}$, invariance of $h_\triangle$ under $z\mapsto\omega z$
implies that its second complex derivative transforms by $\omega^{-2}$;
the same is true of $\wp$.  Their constant difference must therefore vanish.
Expanding each summand for $|z|<a$ gives
\[
 \wp(z)=z^{-2}+\sum_{j\ge1}(j+1)z^j
                         \sum_{\lambda\ne0}\lambda^{-j-2}.
\]
Rotation makes the inner sum zero unless $j+2$ is a multiple of $6$.
Twice complex differentiation of the right side of
\eqref{cal:green-series} gives this series, because
$k_m(k_m-1)c_m=(k_m-1)\sum\lambda^{-k_m}$.
The difference is a real harmonic affine function after subtracting
$\pi|z|^2/2$; rotation removes its linear part and leaves $C_\triangle$.
Normal convergence follows either from this power series or from the
absolutely convergent lattice sums on every smaller disk.
\end{proof}

\subsection{Matching the regular hexagon}
Write $\gamma(r)=-\log r+\pi r^2/2$ for $r>0$.
We first choose the radial data on the regular hexagon, where the test
potential must equal $h_\triangle-C_\triangle$.  In a sector whose
outward normal is the positive real axis, the edge has distance $p$
and the endpoint angles are $\pm\pi/6$.  At an edge point of radius
$t=p\sec\varphi$,
\[
 p+i\sqrt{t^2-p^2}=te^{i|\varphi|}.
\]
The boundary value of the $m$th Green-function mode is therefore
$c_mt^{k_m}\cos(k_m\varphi)$.  Expressing it as a function of the
radius gives the polynomial
\begin{equation}\label{cal:edge-polynomial}
 V_m(r)=c_m\sum_{j=0}^{k_m/2}(-1)^j\binom{k_m}{2j}
                          p^{k_m-2j}(r^2-p^2)^j,
 \qquad p\le r\le R_0.
\end{equation}
On either vertex ray, the same mode equals
$(-1)^mc_mr^{k_m}$, since $k_m\pi/6=m\pi$.  This is the radial
baseline we use for that mode inside the hexagon.  The correction in
\eqref{dual:trial} then recovers the angular dependence exactly:
\[
 (-1)^mc_mr^{k_m}
   +(r/t)^{k_m}\bigl(V_m(t)-(-1)^mc_mt^{k_m}\bigr)
       =c_mr^{k_m}\cos(k_m\varphi).
\]
Thus the edge data $b_m=V_m$ and the vertex-ray baseline determine the
desired potential on each regular sector.  Its remaining part $\gamma$
is already radial and needs no correction there.

\subsection{Continuation to arbitrary sectors}
Other Voronoi sectors can have edge radii below $p$ or above $R_0$.
We now extend the radial data to those radii, keeping two derivatives
and making the correction vanish beyond a fixed radius.  The assembled
potential is required to agree with $h_\triangle-C_\triangle$ only
inside the regular hexagon; the continuation is chosen for the scalar
estimates, not as an extension of that Green function.

Define the clipped polynomial
\[
 S(y)=\begin{cases}
 0,&y\le0,\\
 10y^3-15y^4+6y^5,&0<y<1,\\
 1,&y\ge1.
 \end{cases}
\]
Let $h_s=3/40$ and $c=3/50$.  For $0\le z\le h_s$ put
\[
 P(z)=z-\frac{6z^3}{h_s^2}+\frac{8z^4}{h_s^3}-\frac{3z^5}{h_s^4},\qquad
 Q(z)=\frac{z^2}{2}-\frac{3z^3}{2h_s}
                       +\frac{3z^4}{2h_s^2}-\frac{z^5}{2h_s^3}.
\]
For a function $V$ with two left derivatives at $R_0$, its continuation
is defined for $r\ge R_0$ by
\begin{equation}\label{cal:continuation}
 (\mathcal HV)(r)=
 \begin{cases}
 V(R_0)+P(r-R_0)V'(R_0)+Q(r-R_0)V''(R_0),&r\le R_0+h_s,\\
 V(R_0),&r>R_0+h_s.
 \end{cases}
\end{equation}
Indeed $(P,P',P'')(0)=(0,1,0)$ and
$(Q,Q',Q'')(0)=(0,0,1)$; all six values at $h_s$ vanish.
Thus this operation preserves two derivatives at $R_0$ and joins a
constant with two derivatives at $R_0+h_s$.

Define $b_0=\gamma$ up to $R_0$, continued by
\eqref{cal:continuation}.  For $m\ge1$, set $b_m=V_m$ on $[p,R_0]$
and use $\mathcal HV_m$ above $R_0$.
Below $p$ put
\begin{equation}\label{cal:inner-continuation}
 b_m(r)=\begin{cases}
 0,&0<r\le p-c,\\
 S((r-p+c)/c)\bigl(T_{0m}+T_{1m}(r-p)+T_{2m}(r-p)^2\bigr),&p-c<r<p,
 \end{cases}
\end{equation}
where
\[
 \begin{split}
 T_{0m}&=c_mp^{k_m},\\
 T_{1m}&=-\frac{k_m(k_m-1)}pT_{0m},\\
 T_{2m}&=\frac{k_m(k_m-1)}{p^2}
             \left(\frac{(k_m-2)(k_m-3)}6-\frac12\right)T_{0m}.
 \end{split}
\]
Direct differentiation of \eqref{cal:edge-polynomial} gives
$T_{0m}=V_m(p)$, $T_{1m}=V_m'(p)$, and $2T_{2m}=V_m''(p)$.
Consequently every $b_i$ is $C^2$ away from the origin.

Now define
\[
 U_m(r)=(-1)^mc_mr^{k_m}\bigl(1-S((r-R_0)/(1/10))\bigr),\qquad
 \tau(r)=S\left(\frac{r-R_0-1/1000}{3/40}\right),
\]
and set
\begin{equation}\label{cal:radial-data}
 \begin{split}
 u_0&=(1-\tau)\gamma+\tau b_0,\qquad
 u_m=(1-\tau)U_m+\tau b_m\quad(m\ge1),\\
 B&=\sum_{i\ge0}u_i,\qquad D_i=b_i-u_i,\qquad
 l=R_0+\frac{19}{250}.
 \end{split}
\end{equation}
Here the auxiliary mode $0$ lets the radial part $\gamma$ join its
constant continuation by the same correction mechanism as the other
modes.  Its exponent $k_0=6$ makes this correction vanish at the pole
without changing the singular normalization.  The coefficient
estimates proved below justify the sums and their
first three derivatives on every open spline interval.  In particular,
$B,D_i$ are $C^2$ and piecewise $C^3$ away from zero.  Near zero,
$B=\gamma+\sum_{m\ge1}(-1)^mc_mr^{k_m}$; normal convergence of
\eqref{cal:green-series} therefore gives
\begin{equation}\label{cal:radial-endpoints}
 \begin{gathered}
 B(r)=-\log r+O(r^2),\qquad B'(r)=-1/r+O(r)\quad(r\downarrow0),\\
 D_i(r)=0,\qquad B(r)=B(l)\quad(r\ge l).
 \end{gathered}
\end{equation}
The fixed radial data also satisfy \eqref{dual:radial-sum} for every
$d_0>0$.  Indeed, there are only finitely many retained modes, while
for the omitted modes \eqref{tail:mode-derivative-bound} and the
vanishing of $D_i$ beyond $l$ bound the summand on $[d_0,T_0]$ by
$(1+k_i+T_0)w_{k_i}$.  The explicit geometric terms in
\eqref{tail:w-definition} make this series summable.  This conclusion
uses the analytic mode bounds, not the later finite-core estimates.

For a sector with edge distance $d$ and endpoint angles $-\alpha<\beta$,
write $t(\varphi)=d\sec\varphi$ and use the test potential
\begin{equation}\label{cal:trial-potential}
 H(r,\varphi)=B(r)+\sum_{i\ge0}
                   \left(\frac r{t(\varphi)}\right)^{k_i}D_i(t(\varphi)),
 \qquad 0<r<t(\varphi),\quad -\alpha<\varphi<\beta.
\end{equation}
At the edge its trace is $\sum_i b_i(t)$; at a ray ending at a vertex of
radius $r_v$ its trace is $B(r)+\sum_i(r/r_v)^{k_i}D_i(r_v)$.
These expressions explain both trace matching conditions in the glued
Voronoi construction.

\subsection{The sector integral}
The following formulas separate radial integration from sector geometry.
For $t>0$ define
\begin{equation}\label{cal:radial-integrals}
 \begin{split}
 B_1(t)&=B'(t)+1/t,\qquad S_1(t)=\int_0^t B(r)r\,dr,\\
 S_2(t)&=\int_0^t\bigl(-2B_1(r)+rB_1(r)^2\bigr)\,dr,\\
 H_k(t)&=t^{-k}\int_0^tB_1(r)r^k\,dr,\qquad
 Q_i(t)=tD_i'(t)-k_iD_i(t).
 \end{split}
\end{equation}
Here $H_k$ is indexed by the exponent, so the two modes with exponent
six use the same $H_6$.  They remain separate in every sum involving $D_i$.
Put
\begin{equation}\label{cal:f0-e}
 \begin{split}
 f_0(t)={}&-2\pi\left(S_1(t)+t^2\sum_i\frac{D_i(t)}{k_i+2}\right)\\
 &-\frac12\Biggl(\log t+S_2(t)
   +\sum_i\bigl(-2D_i(t)+2k_iD_i(t)H_{k_i}(t)\bigr)\\
 &\qquad+\sum_{i,j}\frac{k_ik_j}{k_i+k_j}D_i(t)D_j(t)\Biggr),\\
 e(t)={}&\frac12\sum_{i,j}\frac{Q_i(t)Q_j(t)}{k_i+k_j}.
 \end{split}
\end{equation}
All sums here and below include index $0$.
The angular coefficient is nonnegative, since
\begin{equation}\label{cal:positive-kernel}
 e(t)=\frac12\int_0^1
          \left(\sum_iQ_i(t)v^{k_i}\right)^2\frac{dv}{v}\ge0.
\end{equation}

Let $J$ denote the sector contribution to the dual energy,
\[
 J=-2\pi\int_{\rm sector}H
            -\frac12\int_{\rm sector}^{\rm ren}|\nabla H|^2,
\]
where the renormalized integral adds $(\alpha+\beta)\log\eta$ before
letting the radial exclusion $\eta$ tend to zero.  Then
\begin{equation}\label{cal:angle-density}
 J=\int_{-\alpha}^{\beta}
          \bigl(f_0(d\sec\varphi)-\tan^2\varphi\,e(d\sec\varphi)\bigr)
                                                                  \,d\varphi.
\end{equation}
To verify this identity, differentiate \eqref{cal:trial-potential}:
\[
 H_r=-\frac1r+B_1(r)+\sum_i\frac{k_i}{r}(r/t)^{k_i}D_i(t),\qquad
 H_\varphi=\tan\varphi\sum_i(r/t)^{k_i}Q_i(t).
\]
The source integral gives $S_1+t^2\sum_iD_i/(k_i+2)$.
The radial square contributes $\log t+S_2$, its cross terms contribute
$-2D_i+2k_iD_iH_{k_i}$, and its quadratic terms have coefficient
$k_ik_j/(k_i+k_j)$.  Integrating the angular square gives
$\tan^2\varphi\sum Q_iQ_j/(k_i+k_j)$.  These are precisely
\eqref{cal:f0-e}.  The estimates below permit termwise integration.

\subsection{Subtracting area and angle}
Given real constants $\lambda,\mu$, define
\begin{equation}\label{cal:primitive}
 \begin{split}
 f(t)&=f_0(t)-\lambda t^2/2-\mu,\qquad t=\sqrt{d^2+s^2},\\
 G(d,s)&=\frac d{t^2}f(t)-\frac{s^2}{dt^2}e(t),\qquad
 g_d(x)=\int_0^xG(d,s)\,ds\quad(d>0,\ x\in\mathbb R).
 \end{split}
\end{equation}
The evenness of $G$ in $s$ makes $g_d$ odd.  If $A$ is the sector area
and $\theta=\alpha+\beta$, substitution $s=d\tan\varphi$ in
\eqref{cal:angle-density} gives
\begin{equation}\label{cal:sector-primitive}
 J-\lambda A-\mu\theta
       =g_d(d\tan\alpha)+g_d(d\tan\beta).
\end{equation}
This identity holds for signed endpoint angles; only
$|\alpha|,|\beta|<\pi/2$ and $\alpha+\beta>0$ are required.

We choose $\lambda,\mu$ so that the adjusted primitive is stationary
at the regular sector $(d,x)=(p,x_0)$, and set $K$ to twice its value
there.  The exact choices are
\begin{equation}\label{cal:constants}
 \begin{split}
 \lambda=\frac1{px_0}\biggl[&\frac p{x_0}\bigl(f_0(R_0)-f_0(p)\bigr)
                 -\frac{x_0}{p}e(R_0)\\
 &+p\int_0^{x_0}\frac{f_0(\sqrt{p^2+s^2})-f_0(p)}{s^2}\,ds
                 +\frac1p\int_0^{x_0}e(\sqrt{p^2+s^2})\,ds\biggr],\\
 \mu={}&f_0(R_0)-\frac13e(R_0)-\lambda R_0^2/2,\qquad
 K=2g_p(x_0).
 \end{split}
\end{equation}
The first integrand extends continuously to $s=0$ with value
$f_0'(p)/(2p)$.

\begin{lemma}\label{cal:stationarity}
For the constants \eqref{cal:constants},
$\nabla_{d,x}g_d(x)=0$ at $(p,x_0)$.
\end{lemma}
\begin{proof}
Let $g_d^0$ denote the primitive before subtracting area and angle.
For fixed $u=\pi/6$, differentiation in the angle representation gives
\[
 \left.\frac d{dd}g_d^0(d\tan u)\right|_{d=p}
 =\int_0^{x_0}\frac{f_0'(t)-s^2e'(t)/p^2}{t}\,ds,
 \qquad t=\sqrt{p^2+s^2}.
\]
Set $V(s)=f_0(t)-f_0(p)$.  Since $V(s)=O(s^2)$, integration by parts
has no boundary term at zero and yields
\[
 \int_0^{x_0}\frac{f_0'(t)}t\,ds
       =\frac{f_0(R_0)-f_0(p)}{x_0}
                         +\int_0^{x_0}\frac{V(s)}{s^2}\,ds,
\]
whereas
\[
 \int_0^{x_0}\frac{s^2e'(t)}{p^2t}\,ds
       =\frac{x_0e(R_0)}{p^2}-\frac1{p^2}\int_0^{x_0}e(t)\,ds.
\]
The difference is the bracket in \eqref{cal:constants} divided by $p$,
hence equals $\lambda x_0$.
The subtracted area along $x=d\tan u$ is $\lambda d^2\tan u/2$;
its derivative at $p$ is also $\lambda x_0$, and the subtracted angle is
constant.  Thus the derivative of $g_d(d\tan u)$ vanishes there.
Finally $\partial_xg_p(x_0)=G(p,x_0)=0$ by the definition of $\mu$,
because $x_0^2/p^2=1/3$.  The fixed-angle derivative is
$\partial_d+(x_0/p)\partial_x$, so $\partial_d$ vanishes as well.
\end{proof}

\begin{lemma}\label{cal:energy-identity}
The calibrated constants satisfy
\begin{equation}\label{cal:triangular-energy}
 W(E_\triangle)=6K+\lambda+2\pi\mu.
\end{equation}
\end{lemma}
\begin{proof}
Take $\mathbb T=\R^2/\Lambda_\triangle$ and $N=1$.  Its regular
Voronoi hexagon supplies six compatible sectors with $d=p$ and
$\alpha=\beta=\pi/6$.  The radial normalization and summability
proved above supply the hypotheses of Proposition~\ref{dual:identity}
with $d_0=p$, and the torus has area one.  This application uses no
square-torus minimizer or incidence theorem.

On each sector $p\le t\le R_0$,
$u_0(r)=\gamma(r)$ and $u_m(r)=(-1)^mc_mr^{k_m}$ for $r\le t$.
The edge trace and mode cancellation established above therefore give
\[
 H(r,\varphi)=\gamma(r)+\sum_{m\ge1}c_mr^{k_m}\cos(k_m\varphi)
             =h_\triangle-C_\triangle.
\]
All six normals differ by multiples of $\pi/3$, which leaves these
cosines unchanged.  Hence $\nabla H=\nabla h_\triangle$ and the
nonnegative square in \eqref{dual:energy-bound} is zero.
Each sector has area $1/6$ and angle $\pi/3$; its right side in
\eqref{cal:sector-primitive} is $2g_p(x_0)=K$.
Summing the six contributions gives the cell energy
$6K+\lambda+2\pi\mu$.  Lemma~\ref{red:periodic-average} identifies
this cell energy with $W(E_\triangle)$, proving
\eqref{cal:triangular-energy}.
\end{proof}

\subsection{The two scalar estimates needed by the geometry}
The remaining analytic task is to establish
\begin{equation}\label{cal:scalar-bounds}
 \begin{gathered}
 K<0,\qquad 2g_d(x)\ge K\quad(d\ge b,\ x\ge0),\qquad b=\frac{28209}{100000},\\
 \int_0^\infty\max\{0,-G(d,s)\}\,ds\le-K\quad(d\ge b).
 \end{gathered}
\end{equation}
The rational endpoint $b=28209/100000$ is below the geometric
threshold $d_*=1/(2\sqrt\pi)$. Thus the domain $d\ge b$ covers all
edge distances supplied by the separated-minimizer geometry; the
elementary comparison is justified in Section~\ref{sec:conclusion}.

The second estimate is essential when the perpendicular from a point
to a Voronoi edge falls outside that edge.  Its role is made precise
in the following implication, which completes the geometric part of
the argument once \eqref{cal:scalar-bounds} has been certified.

\begin{proposition}\label{cal:scalar-to-torus}
Assume \eqref{cal:scalar-bounds} for the functions and constants just
constructed.  Let $\mathbb T=\R^2/\Gamma$ be a flat torus of area $N$,
where $N\ge1$ is an integer, and let $v_1,\ldots,v_N$ be distinct
sites.  Suppose their full periodic lift gives compatible periodic
Voronoi data with lower edge distance $d_0=b$ and at most $6N$
face-edge sectors.  Let $h$ be the mean-zero potential
\eqref{dual:torus-potential}, with energy
$\mathcal W_{\mathbb T}(h)$ defined by \eqref{dual:torus-energy}.
Then
\[
 \mathcal W_{\mathbb T}(h)\ge NW(E_\triangle).
\]
\end{proposition}
\begin{proof}
The radial normalization and compact summability established above
supply the radial hypotheses of Proposition~\ref{dual:assembly}
with $d_0=b$.  Evenness of $G$ and \eqref{cal:scalar-bounds} supply
the hypotheses of Lemma~\ref{dual:signed-endpoints} at this same
lower distance.  Combining its conclusion with
\eqref{cal:sector-primitive} gives
\[
 J(S)\ge K+\lambda A(S)+\mu\theta(S)
\]
for every sector with $d\ge b$, including the signed endpoint cases.
The compatible-torus conclusion \eqref{dual:compatible-bound} now gives
\[
 \mathcal W_{\mathbb T}(h)\ge N(6K+\lambda+2\pi\mu)
                         =NW(E_\triangle).
\]
The equality is the exact triangular calibration
\eqref{cal:triangular-energy}, not a numerical estimate of the
triangular energy.
\end{proof}
